\chapter{The Fuzzy Distributor Equipment and the Fuzzy Way-Below Relation}
\label{ch:fuzzy-distributors-way-below}

In the previous chapter we constructed the fuzzy ideal monad
\[
    T_{\mathbb V}X
    =
    \operatorname{Idl}_{\mathbb V}(X)
\]
from the upper fuzzy character adjunction. Its strict algebras were called fuzzy
domains. Continuous fuzzy domains were defined by the existence of an enriched
left adjoint
\[
    \beta_X:X\to \operatorname{Idl}_{\mathbb V}(X)
\]
to the ideal-algebra map
\[
    \alpha_X:\operatorname{Idl}_{\mathbb V}(X)\to X.
\]
Thus
\[
    \beta_X\dashv \alpha_X.
\]

The basic algebra and continuity data are summarized by the diagrams
\[
\begin{tikzcd}[column sep=huge,row sep=large]
X
    \ar[r,"\eta_X"]
    \ar[dr,"1_X"']
&
\operatorname{Idl}_{\mathbb V}(X)
    \ar[d,"\alpha_X"]
\\
&
X
\end{tikzcd}
\qquad
\begin{tikzcd}[column sep=huge,row sep=large]
X
    \ar[r,bend left=35,"\beta_X"]
&
\operatorname{Idl}_{\mathbb V}(X)
    \ar[l,bend left=35,"\alpha_X"]
\end{tikzcd}
\]
where the first triangle expresses the strict unit law
\[
    \alpha_X\eta_X=1_X.
\]
In the posetal specialization this gave the fuzzy way-below degree
\[
    \operatorname{wb}_X(a,b)=\beta_X(b)(a).
\]

The purpose of this chapter is to explain the equipment-theoretic origin of this
formula. The way-below relation is not added as further structure. It is
constructed from the ideal-algebra map and the principal-ideal embedding by a
right lifting in the fuzzy distributor equipment.

For a fuzzy domain \(X\), the available maps are
\[
    \eta_X:X\to \operatorname{Idl}_{\mathbb V}(X),
    \qquad
    \eta_X(x)=X(-,x),
\]
and
\[
    \alpha_X:\operatorname{Idl}_{\mathbb V}(X)\to X.
\]
The conjoint of \(\alpha_X\) and the companion of \(\eta_X\) are distributors
with the same source and the same target,
\[
    X\nrightarrow \operatorname{Idl}_{\mathbb V}(X).
\]
They fit into the following equipment diagram:
\[
\begin{tikzcd}[column sep=huge,row sep=large]
&
\operatorname{Idl}_{\mathbb V}(X)
\\
X
    \ar[ur,dashed,bend left=13,"\alpha_X^\ast"]
    \ar[ur,dashed,bend right=13,"{(\eta_X)_\ast}"']
\end{tikzcd}
\]
where
\[
    \alpha_X^\ast(b,I)=X(b,\alpha_X I),
    \qquad
    (\eta_X)_\ast(a,I)
    =
    \operatorname{Idl}_{\mathbb V}(X)(\eta_X a,I).
\]

The fuzzy way-below distributor is defined by the right lifting
\[
    \ll_X
    \coloneqq
    [\alpha_X^\ast,(\eta_X)_\ast]
    :
    X\nrightarrow X.
\]
Equivalently, in the enriched distributor bicategory it represents the object of
distributors \(\Theta:X\nrightarrow X\) equipped with a comparison
\[
    \alpha_X^\ast\circ\Theta\to(\eta_X)_\ast.
\]
In the locally posetal reflection, it is the largest distributor \(\Theta\) for
which
\[
    \alpha_X^\ast\circ\Theta\leq(\eta_X)_\ast.
\]

Diagrammatically, a comparison
\[
    \alpha_X^\ast\circ\Theta\to(\eta_X)_\ast
\]
has the form
\[
\begin{tikzcd}[column sep=huge,row sep=large]
X
    \ar[r,dashed,"\Theta"]
    \ar[rr,bend right=18,dashed,"{(\eta_X)_\ast}"']
&
X
    \ar[r,dashed,"\alpha_X^\ast"]
&
\operatorname{Idl}_{\mathbb V}(X)
    \ar[phantom,from=1-2,to=1-3,"\Downarrow" description]
\end{tikzcd}
\]
and the universal property of \(\ll_X\) says that such comparisons are
represented by comparisons
\[
    \Theta\to \ll_X.
\]

Unwinding the right lifting gives
\[
    \ll_X(a,b)
    \cong
    \int_{I\in \operatorname{Idl}_{\mathbb V}(X)}
    \left[
        X(b,\alpha_X I),
        \operatorname{Idl}_{\mathbb V}(X)(\eta_X a,I)
    \right].
\]
By enriched Yoneda,
\[
    \operatorname{Idl}_{\mathbb V}(X)(\eta_X a,I)
    \cong
    I(a),
\]
and hence
\[
    \ll_X(a,b)
    \cong
    \int_{I\in \operatorname{Idl}_{\mathbb V}(X)}
    \left[
        X(b,\alpha_X I),
        I(a)
    \right].
\]
In the posetal reflection this becomes
\[
    \ll_X(a,b)
    =
    \bigmeet_{I\in\operatorname{Idl}_\Lambda(X)}
    \left(
        X(b,\alpha_X I)\multimap I(a)
    \right).
\]

If \(X\) is continuous, with
\[
    \beta_X\dashv \alpha_X,
\]
then this right-lifted distributor is represented by the approximant map:
\[
    \ll_X(a,b)
    \cong
    \operatorname{Idl}_{\mathbb V}(X)(\eta_X a,\beta_X b)
    \cong
    \beta_X(b)(a).
\]
Thus the way-below degree of the previous chapter is the value of a right-lifted
fuzzy distributor.

Throughout this chapter
\[
    (\mathbb V,\tensor,I,[-,-])
\]
is the fuzzy cosmos fixed in the previous chapters. We write
\[
    \underline{\mathbb V}
\]
for its self-enriched category, and we use the same universe conventions as
before: all ends and coends are formed in a universe in which their indexing
fuzzy category is small and the corresponding limits or colimits in
\(\mathbb V\) exist.

All distributor categories below are understood in this universe-bounded sense.
Thus, when \(\operatorname{Idl}_{\mathbb V}(X)\) is not small in the original
universe, we regard it as small in the next universe and form the corresponding
distributor homs, ends, and coends there.

We write
\[
    \Phi:X\nrightarrow Y
\]
for a fuzzy distributor from \(X\) to \(Y\). The arrow \(\nrightarrow\) denotes a
proarrow, not a negated function arrow.

\section{Inherited data from the fuzzy ideal monad}
\label{sec:inherited-data-from-ideal-monad}

This section records the data from the previous chapter that will be used below.
No new structure is imposed on fuzzy domains in this chapter.

\subsection{The enriched construction}
\label{subsec:inherited-data-enriched}

For a small fuzzy category \(X\), the fuzzy ideal object is the full
\(\mathbb V\)-subcategory
\[
    \operatorname{Idl}_{\mathbb V}(X)
    \subseteq
    [X^{\op},\underline{\mathbb V}]
\]
spanned by finite-flat weights. The unit of the fuzzy ideal monad is the
principal-weight embedding
\[
    \eta_X:X\to \operatorname{Idl}_{\mathbb V}(X),
    \qquad
    \eta_X(x)=X(-,x).
\]
This is the restricted Yoneda embedding:
\[
\begin{tikzcd}[column sep=huge,row sep=large]
X
    \ar[rr,"y_X"]
    \ar[dr,swap,"\eta_X"]
&&
{[X^{\op},\underline{\mathbb V}]}
\\
&
{\operatorname{Idl}_{\mathbb V}(X)}
    \ar[from=2-2,to=1-3,hook,"i_X"']
\end{tikzcd}
\]

The multiplication is flattening:
\[
    \mu_X:
    \operatorname{Idl}_{\mathbb V}(\operatorname{Idl}_{\mathbb V}(X))
    \to
    \operatorname{Idl}_{\mathbb V}(X),
\]
with value
\[
    \mu_X(\mathcal I)(x)
    =
    \int^{I\in \operatorname{Idl}_{\mathbb V}(X)}
        \mathcal I(I)\tensor I(x).
\]

The two unit laws are
\[
    \mu_X\eta_{\operatorname{Idl}_{\mathbb V}(X)}
    =
    1_{\operatorname{Idl}_{\mathbb V}(X)},
    \qquad
    \mu_X\operatorname{Idl}_{\mathbb V}(\eta_X)
    =
    1_{\operatorname{Idl}_{\mathbb V}(X)}.
\]
The first of these is displayed by
\[
\begin{tikzcd}[column sep=huge,row sep=large]
\operatorname{Idl}_{\mathbb V}(X)
    \ar[r,"\eta_{\operatorname{Idl}_{\mathbb V}(X)}"]
    \ar[dr,"1_{\operatorname{Idl}_{\mathbb V}(X)}"']
&
\operatorname{Idl}_{\mathbb V}(\operatorname{Idl}_{\mathbb V}(X))
    \ar[d,"\mu_X"]
\\
&
\operatorname{Idl}_{\mathbb V}(X).
\end{tikzcd}
\]
The second is displayed by
\[
\begin{tikzcd}[column sep=huge,row sep=large]
\operatorname{Idl}_{\mathbb V}(X)
    \ar[r,"\operatorname{Idl}_{\mathbb V}(\eta_X)"]
    \ar[dr,"1_{\operatorname{Idl}_{\mathbb V}(X)}"']
&
\operatorname{Idl}_{\mathbb V}(\operatorname{Idl}_{\mathbb V}(X))
    \ar[d,"\mu_X"]
\\
&
\operatorname{Idl}_{\mathbb V}(X).
\end{tikzcd}
\]
The associativity law is
\[
    \mu_X\operatorname{Idl}_{\mathbb V}(\mu_X)
    =
    \mu_X\mu_{\operatorname{Idl}_{\mathbb V}(X)},
\]
displayed by
\[
\begin{tikzcd}[column sep=huge,row sep=large]
\operatorname{Idl}_{\mathbb V}
    (\operatorname{Idl}_{\mathbb V}
    (\operatorname{Idl}_{\mathbb V}(X)))
    \ar[r,"\operatorname{Idl}_{\mathbb V}(\mu_X)"]
    \ar[d,"\mu_{\operatorname{Idl}_{\mathbb V}(X)}"']
&
\operatorname{Idl}_{\mathbb V}
    (\operatorname{Idl}_{\mathbb V}(X))
    \ar[d,"\mu_X"]
\\
\operatorname{Idl}_{\mathbb V}
    (\operatorname{Idl}_{\mathbb V}(X))
    \ar[r,"\mu_X"']
&
\operatorname{Idl}_{\mathbb V}(X).
\end{tikzcd}
\]

A fuzzy domain is a strict algebra for this monad. Thus it is a small fuzzy
category \(X\) equipped with a fuzzy functor
\[
    \alpha_X:
    \operatorname{Idl}_{\mathbb V}(X)
    \to
    X
\]
such that
\[
    \alpha_X\eta_X=1_X
\]
and
\[
    \alpha_X\mu_X
    =
    \alpha_X\operatorname{Idl}_{\mathbb V}(\alpha_X).
\]
The algebra laws are encoded by the diagrams
\[
\begin{tikzcd}[column sep=huge,row sep=large]
X
    \ar[r,"\eta_X"]
    \ar[dr,"1_X"']
&
\operatorname{Idl}_{\mathbb V}(X)
    \ar[d,"\alpha_X"]
\\
&
X
\end{tikzcd}
\]
and
\[
\begin{tikzcd}[column sep=huge,row sep=large]
\operatorname{Idl}_{\mathbb V}
    (\operatorname{Idl}_{\mathbb V}(X))
    \ar[r,"\mu_X"]
    \ar[d,"\operatorname{Idl}_{\mathbb V}(\alpha_X)"']
&
\operatorname{Idl}_{\mathbb V}(X)
    \ar[d,"\alpha_X"]
\\
\operatorname{Idl}_{\mathbb V}(X)
    \ar[r,"\alpha_X"']
&
X.
\end{tikzcd}
\]
The map \(\alpha_X\) evaluates a finite-flat weight by its chosen ideal colimit.

A fuzzy domain \(X\) is continuous when
\[
    \alpha_X:
    \operatorname{Idl}_{\mathbb V}(X)
    \to
    X
\]
has an enriched left adjoint
\[
    \beta_X:X\to \operatorname{Idl}_{\mathbb V}(X),
    \qquad
    \beta_X\dashv\alpha_X.
\]
The object \(\beta_X(b)\) is the finite-flat weight of approximants of \(b\). The
adjunction supplies unit and counit maps
\[
    1_X\to \alpha_X\beta_X,
    \qquad
    \beta_X\alpha_X\to 1_{\operatorname{Idl}_{\mathbb V}(X)}.
\]
These may be represented by
\[
\begin{tikzcd}[column sep=huge,row sep=large]
X
    \ar[r,"\beta_X"]
    \ar[dr,"1_X"']
&
\operatorname{Idl}_{\mathbb V}(X)
    \ar[d,"\alpha_X"]
\\
&
X
\ar[ul,phantom,"\scriptstyle\eta^{\beta\dashv\alpha}" description]
\end{tikzcd}
\qquad
\begin{tikzcd}[column sep=huge,row sep=large]
\operatorname{Idl}_{\mathbb V}(X)
    \ar[r,"\alpha_X"]
    \ar[dr,"1_{\operatorname{Idl}_{\mathbb V}(X)}"']
&
X
    \ar[d,"\beta_X"]
\\
&
\operatorname{Idl}_{\mathbb V}(X)
\ar[ul,phantom,"\scriptstyle\varepsilon^{\beta\dashv\alpha}" description]
\end{tikzcd}
\]

\subsection{The posetal reflection}
\label{subsec:inherited-data-posetal}

Let
\[
    \Lambda=(L,\leq,\meet,\join,\tensor,\multimap,\bot_L,\top_L)
\]
be a fuzzy truth-value object. For a fuzzy preorder \(X\), the object
\[
    \operatorname{Idl}_\Lambda(X)
\]
is a full subpreorder of
\[
    [X^{\op},\Lambda].
\]
Its elements are the finite-flat fuzzy ideals
\[
    I:X^{\op}\to \Lambda.
\]
The principal ideal at \(x\) is
\[
    \eta_X(x)(y)=X(y,x).
\]
A fuzzy domain over \(\Lambda\) has an algebra map
\[
    \alpha_X:
    \operatorname{Idl}_\Lambda(X)
    \to
    X
\]
satisfying the strict algebra laws
\[
    \alpha_X\eta_X=1_X
\]
and
\[
    \alpha_X\mu_X
    =
    \alpha_X\operatorname{Idl}_\Lambda(\alpha_X).
\]
These are strict algebra laws for the transported ideal monad. In the posetal
specialization they are literal equalities of \(\Lambda\)-functors, equivalently
objectwise equalities compatible with the fuzzy preorder structure.

If \(X\) is continuous, then there is a \(\Lambda\)-adjunction
\[
    \beta_X\dashv\alpha_X,
\]
and the previous chapter defined the fuzzy way-below degree by
\[
    \operatorname{wb}_X(a,b)=\beta_X(b)(a).
\]
The goal of the present chapter is to reconstruct this value as a right lifting.

\subsection{Examples}
\label{subsec:inherited-data-examples}

\begin{example}[The crisp case]
\label{ex:inherited-data-crisp}
For
\[
    \Lambda=2,
\]
fuzzy preorders are ordinary preorders and fuzzy ideals are ordinary ideals,
namely directed lower sets. The principal ideal embedding is
\[
    \eta_X(x)=\downarrow x,
\]
and a strict algebra
\[
    \alpha_X:\operatorname{Idl}(X)\to X
\]
assigns to each ideal its supremum:
\[
    \alpha_X(I)=\bigvee I.
\]
Thus separated strict algebras recover dcpos. If \(X\) is continuous, then
\[
    \beta_X(b)=\{a\in X\mid a\ll b\}.
\]
The crisp data fit the diagram
\[
\begin{tikzcd}[column sep=huge,row sep=large]
X
    \ar[r,"x\mapsto\downarrow x"]
    \ar[dr,"1_X"']
&
\operatorname{Idl}(X)
    \ar[d,"I\mapsto\bigvee I"]
\\
&
X.
\end{tikzcd}
\]
\end{example}

\begin{example}[G\"odel truth values]
\label{ex:inherited-data-godel}
For G\"odel truth values,
\[
    r\tensor s=r\meet s.
\]
A fuzzy ideal is a finite-flat lower predicate
\[
    I:X^{\op}\to [0,1],
\]
and \(\alpha_X(I)\) is the chosen G\"odel-valued ideal colimit of \(I\).
\end{example}

\begin{example}[{\L}ukasiewicz truth values]
\label{ex:inherited-data-lukasiewicz}
For {\L}ukasiewicz truth values,
\[
    r\tensor s=\max(0,r+s-1),
    \qquad
    r\multimap s=\min(1,1-r+s).
\]
Finite-flat fuzzy ideals are lower predicates whose flatness is computed using
this bounded-additive tensor.
\end{example}

\begin{example}[Product truth values]
\label{ex:inherited-data-product}
For the product tensor,
\[
    r\tensor s=rs,
\]
and
\[
    r\multimap s
    =
    \begin{cases}
    1,& r\leq s,\\
    s/r,& r>s.
    \end{cases}
\]
The ideal monad is the finite-flat completion for multiplicative fuzzy weights.
\end{example}

\begin{example}[Lawvere truth values]
\label{ex:inherited-data-lawvere}
For the Lawvere truth-value object, the order is
\[
    a\leq_L b
    \quad\Longleftrightarrow\quad
    a\geq b
\]
in the usual numerical order, the tensor is addition, the top element is \(0\),
and the bottom element is \(\infty\). The residual is
\[
    a\multimap c=c\ominus a.
\]
Fuzzy ideals are finite-flat cost weights. The algebra map
\[
    \alpha_X:\operatorname{Idl}_\Lambda(X)\to X
\]
evaluates such weights by their Lawvere-enriched weighted colimits.
\end{example}

\section{Fuzzy distributors}
\label{sec:fuzzy-distributors}

We now construct the horizontal arrows of the fuzzy distributor equipment.

\subsection{The enriched construction}
\label{subsec:fuzzy-distributors-enriched}

\begin{definition}[Fuzzy distributor]
\label{def:fuzzy-distributor}
Let \(X\) and \(Y\) be small fuzzy categories. A \textbf{fuzzy distributor}
\[
    \Phi:X\nrightarrow Y
\]
is a fuzzy functor
\[
    \Phi:X^{\op}\tensor Y\to \underline{\mathbb V}.
\]
We write its value at \((x,y)\) as
\[
    \Phi(x,y).
\]
\end{definition}

Thus a fuzzy distributor is contravariant in its first variable and covariant in
its second variable. This variance is indicated by the diagram
\[
\begin{tikzcd}[column sep=large,row sep=large]
X^{\op}\tensor Y
    \ar[rr,"\Phi"]
&&
\underline{\mathbb V}.
\end{tikzcd}
\]

\begin{definition}[Composition of fuzzy distributors]
\label{def:fuzzy-distributor-composition}
Let
\[
    \Phi:X\nrightarrow Y,
    \qquad
    \Psi:Y\nrightarrow Z
\]
be fuzzy distributors. Their composite
\[
    \Psi\circ\Phi:X\nrightarrow Z
\]
is defined by
\[
    (\Psi\circ\Phi)(x,z)
    =
    \int^{y\in Y}
    \Phi(x,y)\tensor \Psi(y,z).
\]
The identity distributor on \(X\) is
\[
    1_X:X\nrightarrow X,
    \qquad
    1_X(x,x')=X(x,x').
\]
\end{definition}

The composite is represented by the coend-convolution diagram
\[
\begin{tikzcd}[column sep=huge,row sep=large]
X
    \ar[r,dashed,"\Phi"]
    \ar[rr,bend right=18,dashed,"\Psi\circ\Phi"']
&
Y
    \ar[r,dashed,"\Psi"]
&
Z.
\end{tikzcd}
\]
The identity distributor is the horizontal shadow of the hom-object:
\[
\begin{tikzcd}[column sep=large,row sep=large]
X
    \ar[r,dashed,"1_X"]
&
X.
\end{tikzcd}
\]

\begin{proposition}[The fuzzy distributor bicategory]
\label{prop:fuzzy-distributor-bicategory}
Small fuzzy categories in a fixed universe, fuzzy distributors, and enriched
distributor transformations form a bicategory
\[
    \operatorname{Dist}_{\mathbb V}.
\]
For small fuzzy categories \(X\) and \(Y\), the hom \(\mathbb V\)-category of
distributors from \(X\) to \(Y\) is
\[
    \operatorname{Dist}_{\mathbb V}(X,Y)
    =
    [X^{\op}\tensor Y,\underline{\mathbb V}].
\]
Composition is the coend convolution of
Definition~\ref{def:fuzzy-distributor-composition}, and identities are the
hom-distributors.
\end{proposition}

\begin{proof}
The coend formula defines a fuzzy functor
\[
    X^{\op}\tensor Z\to \underline{\mathbb V}
\]
by the functoriality of enriched coends and the functoriality of \(\tensor\).
Associativity follows from Fubini for enriched coends together with associativity
of the tensor product. In diagrammatic form, both routes below evaluate to the
same iterated coend:
\[
\begin{tikzcd}[column sep=huge,row sep=large]
X
    \ar[r,dashed,"\Phi"]
    \ar[rr,bend left=22,dashed,"\Psi\circ\Phi"]
    \ar[rrr,bend right=34,dashed,"{(\Xi\circ\Psi)\circ\Phi}"']
    \ar[rrr,bend left=34,dashed,"{\Xi\circ(\Psi\circ\Phi)}"]
&
Y
    \ar[r,dashed,"\Psi"]
    \ar[rr,bend right=22,swap,dashed,"\Xi\circ\Psi"]
&
Z
    \ar[r,dashed,"\Xi"]
&
W.
\end{tikzcd}
\]
More explicitly,
\[
\begin{aligned}
    \Xi\circ(\Psi\circ\Phi)
    &\cong
    \int^{z}
    \left(
        \int^{y}\Phi(-,y)\tensor\Psi(y,z)
    \right)
    \tensor \Xi(z,-)
\\
    &\cong
    \int^{y,z}
        \Phi(-,y)\tensor\Psi(y,z)\tensor\Xi(z,-)
\\
    &\cong
    (\Xi\circ\Psi)\circ\Phi.
\end{aligned}
\]
The unit laws are the enriched co-Yoneda formulas
\[
    \int^{y\in Y}
        Y(y,y')\tensor \Phi(x,y)
    \cong
    \Phi(x,y')
\]
and
\[
    \int^{x'\in X}
        X(x,x')\tensor \Phi(x',y)
    \cong
    \Phi(x,y).
\]
These are expressed by
\[
\begin{tikzcd}[column sep=huge,row sep=large]
X
    \ar[r,dashed,"\Phi"]
    \ar[rr,bend right=18,dashed,"\Phi"']
&
Y
    \ar[r,dashed,"1_Y"]
&
Y
\end{tikzcd}
\qquad
\begin{tikzcd}[column sep=huge,row sep=large]
X
    \ar[r,dashed,"1_X"]
    \ar[rr,bend right=18,dashed,"\Phi"']
&
X
    \ar[r,dashed,"\Phi"]
&
Y.
\end{tikzcd}
\]
\end{proof}

Together with fuzzy functors as vertical arrows and the companion--conjoint
assignments below, this bicategory underlies the fuzzy distributor equipment.

A cell in the equipment
\[
\begin{tikzcd}[column sep=huge,row sep=large]
X
    \ar[d,"f"']
    \ar[r,dashed,"\Phi"]
&
Y
    \ar[d,"g"]
\\
X'
    \ar[r,dashed,"\Psi"']
&
Y'
\ar[phantom,from=1-2,to=2-1,"\Downarrow" description]
\end{tikzcd}
\]
is a distributor transformation
\[
    \Phi(x,y)\to \Psi(fx,gy),
\]
natural in \(x\in X\) and \(y\in Y\). In the locally posetal reflection, this is
the pointwise inequality
\[
    \Phi(x,y)\leq \Psi(fx,gy).
\]

\subsection{The posetal reflection}
\label{subsec:fuzzy-distributors-posetal}

In the posetal specialization determined by
\[
    \Lambda=(L,\leq,\meet,\join,\tensor,\multimap,\bot_L,\top_L),
\]
a fuzzy distributor
\[
    \Phi:X\nrightarrow Y
\]
is a map
\[
    \Phi:X\times Y\to L
\]
satisfying
\[
    X(x',x)\tensor \Phi(x,y)\tensor Y(y,y')
    \leq
    \Phi(x',y')
\]
for all \(x,x'\in X\) and \(y,y'\in Y\). Thus \(\Phi\) is lower in the first
variable and upper in the second variable.

Composition is
\[
    (\Psi\circ\Phi)(x,z)
    =
    \bigjoin_{y\in Y}
    \Phi(x,y)\tensor\Psi(y,z),
\]
and the identity distributor is
\[
    1_X(x,x')=X(x,x').
\]

\subsection{Examples}
\label{subsec:fuzzy-distributors-examples}

\begin{example}[The crisp case]
\label{ex:fuzzy-distributors-crisp}
For \(\Lambda=2\), fuzzy categories are preorders and a distributor
\[
    \Phi:X\nrightarrow Y
\]
is a relation
\[
    \Phi\subseteq X\times Y
\]
such that
\[
    x'\leq_X x,\quad \Phi(x,y),\quad y\leq_Y y'
    \quad\Longrightarrow\quad
    \Phi(x',y').
\]
Composition is ordinary relational composition:
\[
    (\Psi\circ\Phi)(x,z)
    \Longleftrightarrow
    \exists y\in Y,\; \Phi(x,y)\wedge \Psi(y,z).
\]
\end{example}

\begin{example}[G\"odel-valued distributors]
\label{ex:fuzzy-distributors-godel}
For G\"odel truth values,
\[
    (\Psi\circ\Phi)(x,z)
    =
    \bigjoin_{y\in Y}
    \min(\Phi(x,y),\Psi(y,z)).
\]
\end{example}

\begin{example}[{\L}ukasiewicz-valued distributors]
\label{ex:fuzzy-distributors-lukasiewicz}
For {\L}ukasiewicz truth values,
\[
    (\Psi\circ\Phi)(x,z)
    =
    \bigjoin_{y\in Y}
    \max(0,\Phi(x,y)+\Psi(y,z)-1).
\]
\end{example}

\begin{example}[Product-valued distributors]
\label{ex:fuzzy-distributors-product}
For product truth values,
\[
    (\Psi\circ\Phi)(x,z)
    =
    \bigjoin_{y\in Y}
    \Phi(x,y)\Psi(y,z).
\]
\end{example}

\begin{example}[Lawvere-valued distributors]
\label{ex:fuzzy-distributors-lawvere}
For the Lawvere quantale, joins in the truth-value order are infima in the usual
numerical order, and tensor is addition. Hence distributor composition is
infimal convolution:
\[
    (\Psi\circ\Phi)(x,z)
    =
    \inf_{y\in Y}
    \bigl(\Phi(x,y)+\Psi(y,z)\bigr).
\]
\end{example}

\section{Companions and conjoints}
\label{sec:companions-conjoints}

Fuzzy functors determine distinguished fuzzy distributors. These are the
companion and conjoint constructions. They place vertical arrows inside the
horizontal distributor equipment.

\subsection{The enriched construction}
\label{subsec:companions-conjoints-enriched}

\begin{definition}[Companion and conjoint]
\label{def:companion-conjoint}
Let
\[
    f:X\to Y
\]
be a fuzzy functor. Its \textbf{companion} is the distributor
\[
    f_\ast:X\nrightarrow Y
\]
defined by
\[
    f_\ast(x,y)=Y(fx,y).
\]
Its \textbf{conjoint} is the distributor
\[
    f^\ast:Y\nrightarrow X
\]
defined by
\[
    f^\ast(y,x)=Y(y,fx).
\]
\end{definition}

The companion and conjoint sit on opposite sides of \(f\):
\[
\begin{tikzcd}[column sep=huge,row sep=large]
X
    \ar[d,"f"']
    \ar[r,dashed,"f_\ast"]
&
Y
    \ar[d,"1_Y"]
\\
Y
    \ar[r,dashed,"1_Y"']
&
Y
\end{tikzcd}
\qquad
\begin{tikzcd}[column sep=huge,row sep=large]
Y
    \ar[d,"1_Y"']
    \ar[r,dashed,"f^\ast"]
&
X
    \ar[d,"f"]
\\
Y
    \ar[r,dashed,"1_Y"']
&
Y.
\end{tikzcd}
\]
Equivalently,
\[
\begin{tikzcd}[column sep=huge,row sep=large]
X
    \ar[r,bend left=14,dashed,"f_\ast"]
&
Y
    \ar[l,bend left=14,dashed,"f^\ast"]
\end{tikzcd}
\]
is an adjoint pair in the distributor bicategory.

\begin{proposition}[Companion--conjoint adjunction]
\label{prop:companion-conjoint-adjunction}
For every fuzzy functor
\[
    f:X\to Y,
\]
there is an adjunction in the distributor bicategory
\[
    f_\ast\dashv f^\ast.
\]
\end{proposition}

\begin{proof}
The unit is a distributor transformation
\[
    1_X\to f^\ast\circ f_\ast.
\]
At \((x,x')\), the target is
\[
    (f^\ast\circ f_\ast)(x,x')
    =
    \int^{y\in Y}
        Y(fx,y)\tensor Y(y,fx').
\]
The required map is obtained by composing the functoriality map
\[
    X(x,x')\to Y(fx,fx')
\]
with the coend coprojection at \(y=fx'\):
\[
    Y(fx,fx')
    \cong
    Y(fx,fx')\tensor I
    \to
    Y(fx,fx')\tensor Y(fx',fx')
    \to
    \int^{y\in Y}
        Y(fx,y)\tensor Y(y,fx').
\]

The counit is a distributor transformation
\[
    f_\ast\circ f^\ast\to 1_Y.
\]
At \((y,y')\), it is induced by composition in \(Y\):
\[
    \int^{x\in X}
        Y(y,fx)\tensor Y(fx,y')
    \to
    Y(y,y').
\]
The triangle identities follow from associativity and unitality of composition in
\(Y\), equivalently from the standard companion--conjoint construction in the
enriched distributor equipment.
\end{proof}

\begin{proposition}[Functoriality of companions and conjoints]
\label{prop:companion-conjoint-functoriality}
For fuzzy functors
\[
    X\xrightarrow{f}Y\xrightarrow{g}Z,
\]
there are canonical isomorphisms of distributors
\[
    (gf)_\ast
    \cong
    g_\ast\circ f_\ast
\]
and
\[
    (gf)^\ast
    \cong
    f^\ast\circ g^\ast.
\]
Moreover,
\[
    (1_X)_\ast\cong 1_X,
    \qquad
    (1_X)^\ast\cong 1_X.
\]
\end{proposition}

\begin{proof}
For companions,
\[
\begin{aligned}
    (g_\ast\circ f_\ast)(x,z)
    &=
    \int^{y\in Y}
        Y(fx,y)\tensor Z(gy,z)
\\
    &\cong
    Z(gfx,z)
\\
    &=
    (gf)_\ast(x,z),
\end{aligned}
\]
by enriched Yoneda. Diagrammatically:
\[
\begin{tikzcd}[column sep=huge,row sep=large]
X
    \ar[r,dashed,"f_\ast"]
    \ar[rr,bend right=18,dashed,"{(gf)_\ast}"']
&
Y
    \ar[r,dashed,"g_\ast"]
&
Z.
\end{tikzcd}
\]
For conjoints,
\[
\begin{aligned}
    (f^\ast\circ g^\ast)(z,x)
    &=
    \int^{y\in Y}
        Z(z,gy)\tensor Y(y,fx)
\\
    &\cong
    Z(z,gfx)
\\
    &=
    (gf)^\ast(z,x).
\end{aligned}
\]
This is represented by the reverse diagram
\[
\begin{tikzcd}[column sep=huge,row sep=large]
Z
    \ar[r,dashed,"g^\ast"]
    \ar[rr,bend right=18,dashed,"{(gf)^\ast}"']
&
Y
    \ar[r,dashed,"f^\ast"]
&
X.
\end{tikzcd}
\]
The identity statements follow directly from the definitions.
\end{proof}

\subsection{The posetal reflection}
\label{subsec:companions-conjoints-posetal}

For a \(\Lambda\)-functor
\[
    f:X\to Y,
\]
the companion and conjoint are
\[
    f_\ast(x,y)=Y(fx,y)
\]
and
\[
    f^\ast(y,x)=Y(y,fx).
\]
The adjunction
\[
    f_\ast\dashv f^\ast
\]
amounts to the distributor inequalities
\[
    1_X\leq f^\ast\circ f_\ast
\]
and
\[
    f_\ast\circ f^\ast\leq 1_Y.
\]
The first inequality follows from functoriality of \(f\), and the second follows
from transitivity in \(Y\).

\subsection{Examples}
\label{subsec:companions-conjoints-examples}

\begin{example}[The crisp case]
\label{ex:companions-crisp}
For a monotone map
\[
    f:X\to Y
\]
between preorders,
\[
    f_\ast(x,y)
    \Longleftrightarrow
    f(x)\leq_Y y,
\]
and
\[
    f^\ast(y,x)
    \Longleftrightarrow
    y\leq_Y f(x).
\]
\end{example}

\begin{example}[Standard fuzzy truth values]
\label{ex:companions-standard}
For G\"odel, {\L}ukasiewicz, and product truth values, the formulas remain
\[
    f_\ast(x,y)=Y(fx,y),
    \qquad
    f^\ast(y,x)=Y(y,fx).
\]
The tensor and residual of the chosen truth-value object determine the distributor
composition and the adjunction inequalities.
\end{example}

\begin{example}[Lawvere truth values]
\label{ex:companions-lawvere}
For the Lawvere quantale, a fuzzy functor is a nonexpansive map. The companion
and conjoint of
\[
    f:X\to Y
\]
are
\[
    f_\ast(x,y)=d_Y(fx,y),
    \qquad
    f^\ast(y,x)=d_Y(y,fx).
\]
Thus companions and conjoints record the forward and backward distance kernels
induced by \(f\).
\end{example}

\section{Right liftings of fuzzy distributors}
\label{sec:right-liftings-fuzzy-distributors}

Right liftings are the categorical operation that will produce the fuzzy
way-below distributor.

\subsection{The enriched construction}
\label{subsec:right-liftings-enriched}

Let
\[
    \Phi:A\nrightarrow C,
    \qquad
    \Psi:B\nrightarrow C
\]
be fuzzy distributors with common target \(C\).

\begin{definition}[Right lifting]
\label{def:right-lifting}
The \textbf{right lifting} of \(\Psi\) through \(\Phi\) is the fuzzy distributor
\[
    [\Phi,\Psi]:B\nrightarrow A
\]
defined by
\[
    [\Phi,\Psi](b,a)
    =
    \int_{c\in C}
    [\Phi(a,c),\Psi(b,c)].
\]
\end{definition}

The right lifting is characterized by the following universal diagram:
\[
\begin{tikzcd}[column sep=huge,row sep=large]
B
    \ar[r,dashed,"\Theta"]
    \ar[rr,bend right=18,dashed,"\Psi"']
&
A
    \ar[r,dashed,"\Phi"]
&
C
\ar[phantom,from=1-2,to=1-3,"\Downarrow" description]
\end{tikzcd}
\]
where comparisons
\[
    \Theta\to[\Phi,\Psi]
\]
correspond to comparisons
\[
    \Phi\circ\Theta\to\Psi.
\]
In the locally posetal reflection, this says that \([\Phi,\Psi]\) is the largest
\(\Theta\) satisfying
\[
    \Phi\circ\Theta\leq\Psi.
\]

\begin{proposition}[Universal property of right liftings]
\label{prop:right-lifting-universal-property}
For every fuzzy distributor
\[
    \Theta:B\nrightarrow A,
\]
there is a natural isomorphism
\[
    \operatorname{Dist}_{\mathbb V}(B,A)(\Theta,[\Phi,\Psi])
    \cong
    \operatorname{Dist}_{\mathbb V}(B,C)(\Phi\circ\Theta,\Psi).
\]
Equivalently, postcomposition by \(\Phi\) has right adjoint
\[
    [\Phi,-].
\]
\end{proposition}

\begin{proof}
Using the enriched hom-object in the distributor category, the left-hand side is
\[
\begin{aligned}
\operatorname{Dist}_{\mathbb V}(B,A)(\Theta,[\Phi,\Psi])
&=
\int_{b,a}
[
    \Theta(b,a),
    [\Phi,\Psi](b,a)
]
\\
&=
\int_{b,a}
\left[
    \Theta(b,a),
    \int_c[\Phi(a,c),\Psi(b,c)]
\right]
\\
&\cong
\int_{b,a,c}
[
    \Theta(b,a),
    [\Phi(a,c),\Psi(b,c)]
]
\\
&\cong
\int_{b,a,c}
[
    \Theta(b,a)\tensor\Phi(a,c),
    \Psi(b,c)
].
\end{aligned}
\]
By the end--coend adjunction, this is isomorphic to
\[
    \int_{b,c}
    \left[
        \int^{a}
        \Theta(b,a)\tensor\Phi(a,c),
        \Psi(b,c)
    \right].
\]
The coend is
\[
    (\Phi\circ\Theta)(b,c).
\]
Hence the displayed object is
\[
    \operatorname{Dist}_{\mathbb V}(B,C)(\Phi\circ\Theta,\Psi).
\]
\end{proof}

\subsection{The posetal reflection}
\label{subsec:right-liftings-posetal}

In the posetal reflection, the right lifting is
\[
    [\Phi,\Psi](b,a)
    =
    \bigmeet_{c\in C}
    \left(
        \Phi(a,c)\multimap\Psi(b,c)
    \right).
\]
Its universal property becomes
\[
    \Theta\leq[\Phi,\Psi]
    \quad\Longleftrightarrow\quad
    \Phi\circ\Theta\leq\Psi.
\]
Thus \([\Phi,\Psi](b,a)\) is the degree to which every instance of
\(\Phi(a,c)\) entails the corresponding instance of \(\Psi(b,c)\).

The posetal lifting adjunction can be displayed as
\[
\begin{tikzcd}[column sep=huge,row sep=large]
\operatorname{Dist}_{\Lambda}(B,A)
    \ar[r,bend left=35,"{\Phi\circ -}"]
&
\operatorname{Dist}_{\Lambda}(B,C)
    \ar[l,bend left=35,"{[\Phi,-]}"]
\end{tikzcd}
\]
with order-enriched adjunction condition
\[
    \Phi\circ\Theta\leq\Psi
    \quad\Longleftrightarrow\quad
    \Theta\leq[\Phi,\Psi].
\]

\subsection{Examples}
\label{subsec:right-liftings-examples}

\begin{example}[The crisp case]
\label{ex:right-lifting-crisp}
For \(\Lambda=2\),
\[
    [\Phi,\Psi](b,a)
\]
holds precisely when
\[
    \forall c\in C,\;
    \Phi(a,c)\Rightarrow \Psi(b,c).
\]
The universal property is
\[
    \Theta\subseteq[\Phi,\Psi]
    \quad\Longleftrightarrow\quad
    \Phi\circ\Theta\subseteq\Psi.
\]
\end{example}

\begin{example}[G\"odel right liftings]
\label{ex:right-lifting-godel}
For G\"odel truth values,
\[
    [\Phi,\Psi](b,a)
    =
    \bigmeet_{c\in C}
    \left(
        \Phi(a,c)\multimap_{\mathrm G}\Psi(b,c)
    \right),
\]
where
\[
    u\multimap_{\mathrm G}v
    =
    \begin{cases}
    1,&u\leq v,\\
    v,&u>v.
    \end{cases}
\]
\end{example}

\begin{example}[{\L}ukasiewicz right liftings]
\label{ex:right-lifting-lukasiewicz}
For {\L}ukasiewicz truth values,
\[
    [\Phi,\Psi](b,a)
    =
    \bigmeet_{c\in C}
    \min(1,1-\Phi(a,c)+\Psi(b,c)).
\]
\end{example}

\begin{example}[Product right liftings]
\label{ex:right-lifting-product}
For product truth values,
\[
    [\Phi,\Psi](b,a)
    =
    \bigmeet_{c\in C}
    \left(
        \Phi(a,c)\multimap_\Pi\Psi(b,c)
    \right),
\]
where
\[
    u\multimap_\Pi v
    =
    \begin{cases}
    1,&u\leq v,\\
    v/u,&u>v.
    \end{cases}
\]
\end{example}

\begin{example}[Lawvere right liftings]
\label{ex:right-lifting-lawvere}
For the Lawvere quantale, the internal hom is truncated subtraction. Written in
the usual numerical order,
\[
    [\Phi,\Psi](b,a)
    =
    \sup_{c\in C}
    \bigl(\Psi(b,c)\ominus \Phi(a,c)\bigr),
\]
where
\[
    v\ominus u=\max(0,v-u)
\]
with the usual extended-value conventions. Thus a Lawvere right lifting measures
the largest residual excess needed to transform \(\Phi(a,-)\) into
\(\Psi(b,-)\).
\end{example}

\section{Ideals as distributors and approximable kernels}
\label{sec:ideals-as-distributors}

We next relate the ideal monad from the previous chapter to the distributor
equipment.

\subsection{The enriched construction}
\label{subsec:ideals-as-distributors-enriched}

A fuzzy ideal of \(X\) is a finite-flat weight
\[
    I:X^{\op}\to\underline{\mathbb V}.
\]
Equivalently, it is a fuzzy distributor
\[
    I:X\nrightarrow \mathbf 1,
\]
because
\[
    X^{\op}\tensor\mathbf 1
    \cong
    X^{\op}.
\]
Thus
\[
    \operatorname{Idl}_{\mathbb V}(X)
\]
is the full \(\mathbb V\)-subcategory of
\[
    \operatorname{Dist}_{\mathbb V}(X,\mathbf 1)
    =
    [X^{\op},\underline{\mathbb V}]
\]
spanned by finite-flat distributors. This inclusion is represented by
\[
\begin{tikzcd}[column sep=huge,row sep=large]
{\operatorname{Idl}_{\mathbb V}(X)}
    \ar[r,hook,"i_X"]
&
{\operatorname{Dist}_{\mathbb V}(X,\mathbf{1})}
    \ar[r,"\cong"]
&
{[X^{\op},\underline{\mathbb V}]}.
\end{tikzcd}
\]

\begin{definition}[Fuzzy approximable map]
\label{def:fuzzy-approximable-map-equipment}
Let \(X\) and \(Y\) be small fuzzy categories. A \textbf{fuzzy approximable map}
from \(X\) to \(Y\) is a Kleisli morphism for the fuzzy ideal monad:
\[
    R:X\to \operatorname{Idl}_{\mathbb V}(Y).
\]
Equivalently, it is a kernel
\[
    R:X\tensor Y^{\op}\to\underline{\mathbb V}
\]
such that, for every \(x\in X\), the fibre
\[
    R(x,-):Y^{\op}\to\underline{\mathbb V}
\]
is finite-flat.
\end{definition}

The distributor variance of this kernel is opposite to the variance of a
distributor \(X\nrightarrow Y\). Under the convention
\[
    \Phi:A\nrightarrow B
    \quad\text{means}\quad
    \Phi:A^{\op}\tensor B\to\underline{\mathbb V},
\]
the kernel \(R\) corresponds to a distributor
\[
    R^\dagger:Y\nrightarrow X,
    \qquad
    R^\dagger(y,x)=R(x,y).
\]
The variance conversion is summarized by
\[
\begin{tikzcd}[column sep=huge,row sep=large]
X
    \ar[r,"R"]
&
{\operatorname{Idl}_{\mathbb V}(Y)}
    \ar[r,hook,"i_Y"]
&
{[Y^{\op},\underline{\mathbb V}]}
\end{tikzcd}
\qquad
\begin{tikzcd}[column sep=huge,row sep=large]
Y
    \ar[r,dashed,"R^\dagger"]
&
X.
\end{tikzcd}
\]
The upper row is vertical/Kleisli data, while the lower row is distributor data.

\begin{proposition}[Kleisli extension as distributor convolution]
\label{prop:kleisli-extension-convolution}
Let
\[
    R:X\rightsquigarrow Y
\]
be a fuzzy approximable map. Its Kleisli extension
\[
    \overline R:
    \operatorname{Idl}_{\mathbb V}(X)
    \to
    \operatorname{Idl}_{\mathbb V}(Y)
\]
is given by
\[
    \overline R(I)(y)
    =
    \int^{x\in X}
    I(x)\tensor R(x,y).
\]
Equivalently, if
\[
    R^\dagger:Y\nrightarrow X
\]
is the associated oppositely oriented distributor, then
\[
    \overline R(I)
    =
    I\circ R^\dagger.
\]
\end{proposition}

\begin{proof}
The Kleisli extension is
\[
\begin{tikzcd}[column sep=huge,row sep=large]
\operatorname{Idl}_{\mathbb V}(X)
    \ar[r,"\operatorname{Idl}_{\mathbb V}(R)"]
&
\operatorname{Idl}_{\mathbb V}
    (\operatorname{Idl}_{\mathbb V}(Y))
    \ar[r,"\mu_Y"]
&
\operatorname{Idl}_{\mathbb V}(Y).
\end{tikzcd}
\]
The multiplication \(\mu_Y\) is flattening, so
\[
    \overline R(I)(y)
    =
    \int^{x\in X}
    I(x)\tensor R(x,y).
\]
On the other hand,
\[
    (I\circ R^\dagger)(y)
    =
    \int^{x\in X}
    R^\dagger(y,x)\tensor I(x)
    =
    \int^{x\in X}
    R(x,y)\tensor I(x),
\]
which is canonically isomorphic to the same coend by symmetry of \(\tensor\).
\end{proof}

The convolution picture is
\[
\begin{tikzcd}[column sep=huge,row sep=large]
Y
    \ar[r,dashed,"R^\dagger"]
    \ar[rr,bend right=18,dashed,"I\circ R^\dagger"']
&
X
    \ar[r,dashed,"I"]
&
\mathbf 1.
\end{tikzcd}
\]

\subsection{The posetal reflection}
\label{subsec:ideals-as-distributors-posetal}

In the posetal specialization, a fuzzy approximable map
\[
    R:X\rightsquigarrow Y
\]
is a map
\[
    R:X\times Y\to L
\]
such that each fibre
\[
    R(x,-):Y^{\op}\to \Lambda
\]
is a fuzzy ideal of \(Y\), and such that the assignment
\[
    x\longmapsto R(x,-)
\]
is a \(\Lambda\)-functor
\[
    X\to \operatorname{Idl}_\Lambda(Y).
\]
Its extension is
\[
    \overline R(I)(y)
    =
    \bigjoin_{x\in X}
    I(x)\tensor R(x,y).
\]
Kleisli composition is therefore
\[
    (S\circ R)(x,z)
    =
    \bigjoin_{y\in Y}
    R(x,y)\tensor S(y,z).
\]

\subsection{Examples}
\label{subsec:ideals-as-distributors-examples}

\begin{example}[The crisp case]
\label{ex:ideals-as-distributors-crisp}
For \(\Lambda=2\), a fuzzy approximable map
\[
    R:X\rightsquigarrow Y
\]
is a monotone map
\[
    X\to\operatorname{Idl}(Y).
\]
Equivalently, it is a relation
\[
    R\subseteq X\times Y
\]
such that every fibre
\[
    R[x]=\{y\in Y\mid R(x,y)\}
\]
is an ideal of \(Y\), and
\[
    x\leq x'
    \quad\Longrightarrow\quad
    R[x]\subseteq R[x'].
\]
The extension formula becomes
\[
    \overline R(I)=\bigcup_{x\in I}R[x].
\]
\end{example}

\begin{example}[G\"odel approximable maps]
\label{ex:ideals-as-distributors-godel}
For G\"odel truth values,
\[
    \overline R(I)(y)
    =
    \bigjoin_{x\in X}
    \min(I(x),R(x,y)).
\]
\end{example}

\begin{example}[{\L}ukasiewicz approximable maps]
\label{ex:ideals-as-distributors-lukasiewicz}
For {\L}ukasiewicz truth values,
\[
    \overline R(I)(y)
    =
    \bigjoin_{x\in X}
    \max(0,I(x)+R(x,y)-1).
\]
\end{example}

\begin{example}[Product approximable maps]
\label{ex:ideals-as-distributors-product}
For product truth values,
\[
    \overline R(I)(y)
    =
    \bigjoin_{x\in X}
    I(x)R(x,y).
\]
\end{example}

\begin{example}[Lawvere approximable maps]
\label{ex:ideals-as-distributors-lawvere}
For the Lawvere quantale, the extension formula is infimal convolution:
\[
    \overline R(I)(y)
    =
    \inf_{x\in X}
    \bigl(I(x)+R(x,y)\bigr),
\]
written in the usual numerical order.
\end{example}

\section{The ideal algebra map inside the distributor equipment}
\label{sec:ideal-algebra-map-equipment}

We now place the ideal-algebra map and the principal-ideal embedding inside the
distributor equipment.

\subsection{The enriched construction}
\label{subsec:ideal-algebra-map-equipment-enriched}

Let \(X\) be a fuzzy domain. We have the algebra map
\[
    \alpha_X:\operatorname{Idl}_{\mathbb V}(X)\to X
\]
and the principal-ideal embedding
\[
    \eta_X:X\to\operatorname{Idl}_{\mathbb V}(X).
\]

The conjoint of \(\alpha_X\) is the distributor
\[
    \alpha_X^\ast:
    X\nrightarrow \operatorname{Idl}_{\mathbb V}(X)
\]
given by
\[
    \alpha_X^\ast(b,I)
    =
    X(b,\alpha_X I).
\]
The companion of \(\eta_X\) is the distributor
\[
    (\eta_X)_\ast:
    X\nrightarrow \operatorname{Idl}_{\mathbb V}(X)
\]
given by
\[
    (\eta_X)_\ast(a,I)
    =
    \operatorname{Idl}_{\mathbb V}(X)(\eta_X a,I).
\]

The two constructions arise from opposite sides of the algebra triangle:
\[
\begin{tikzcd}[column sep=huge,row sep=large]
X
    \ar[r,"\eta_X"]
    \ar[dr,"1_X"']
&
\operatorname{Idl}_{\mathbb V}(X)
    \ar[d,"\alpha_X"]
\\
&
X
\end{tikzcd}
\qquad
\begin{tikzcd}[column sep=huge,row sep=large]
&
\operatorname{Idl}_{\mathbb V}(X)
\\
X
    \ar[ur,dashed,bend left=14,"\alpha_X^\ast"]
    \ar[ur,dashed,bend right=14,"{(\eta_X)_\ast}"']
\end{tikzcd}
\]

Thus
\[
    \alpha_X^\ast
    \qquad\text{and}\qquad
    (\eta_X)_\ast
\]
are two fuzzy distributors with the same source and the same target:
\[
    X\nrightarrow \operatorname{Idl}_{\mathbb V}(X).
\]
Consequently, the right lifting
\[
    [\alpha_X^\ast,(\eta_X)_\ast]
\]
is a fuzzy distributor
\[
    X\nrightarrow X.
\]
This is the distributor that will be defined as the fuzzy way-below relation.

\subsection{The posetal reflection}
\label{subsec:ideal-algebra-map-equipment-posetal}

For a fuzzy domain over \(\Lambda\), the two distributors are
\[
    \alpha_X^\ast(b,I)
    =
    X(b,\alpha_X I),
\]
and
\[
    (\eta_X)_\ast(a,I)
    =
    \operatorname{Idl}_\Lambda(X)(\eta_X a,I).
\]
By the enriched Yoneda lemma in the posetal setting,
\[
    \operatorname{Idl}_\Lambda(X)(\eta_X a,I)
    =
    I(a).
\]
Hence
\[
    (\eta_X)_\ast(a,I)=I(a).
\]

\subsection{Examples}
\label{subsec:ideal-algebra-map-equipment-examples}

\begin{example}[The crisp case]
\label{ex:ideal-algebra-map-crisp}
For \(\Lambda=2\),
\[
    \alpha_X^\ast(b,I)
    \Longleftrightarrow
    b\leq \bigvee I,
\]
and
\[
    (\eta_X)_\ast(a,I)
    \Longleftrightarrow
    \downarrow a\subseteq I.
\]
Since \(I\) is lower,
\[
    \downarrow a\subseteq I
    \quad\Longleftrightarrow\quad
    a\in I.
\]
Thus the right lifting expresses the ordinary implication
\[
    b\leq\bigvee I
    \quad\Longrightarrow\quad
    a\in I
\]
uniformly over all ideals \(I\).
\end{example}

\begin{example}[Standard fuzzy truth values]
\label{ex:ideal-algebra-map-standard}
For G\"odel, {\L}ukasiewicz, and product truth values, the formulas are
\[
    \alpha_X^\ast(b,I)=X(b,\alpha_X I),
    \qquad
    (\eta_X)_\ast(a,I)=I(a).
\]
The later right-lifting formula compares these two degrees using the residual
appropriate to the chosen tensor.
\end{example}

\begin{example}[Lawvere truth values]
\label{ex:ideal-algebra-map-lawvere}
For the Lawvere quantale,
\[
    \alpha_X^\ast(b,I)=d_X(b,\alpha_X I)
\]
is the cost from \(b\) to the ideal colimit of \(I\), while
\[
    (\eta_X)_\ast(a,I)=I(a)
\]
is the cost-degree with which \(a\) belongs to the weight \(I\).
\end{example}

\section{The fuzzy way-below distributor}
\label{sec:fuzzy-way-below-distributor}

The fuzzy way-below relation is now obtained by applying the right-lifting
construction to the two distributors built from \(\alpha_X\) and \(\eta_X\).

\subsection{The enriched construction}
\label{subsec:fuzzy-way-below-enriched}

\begin{definition}[Fuzzy way-below distributor]
\label{def:fuzzy-way-below-distributor}
Let \(X\) be a fuzzy domain. The \textbf{fuzzy way-below distributor} of \(X\)
is the right lifting
\[
    \ll_X
    \coloneqq
    [\alpha_X^\ast,(\eta_X)_\ast]
    :
    X\nrightarrow X.
\]
We write
\[
    \ll_X(a,b)
\]
for its value at \((a,b)\), where \(a\) is the approximant variable and \(b\) is
the element being approximated.
\end{definition}

With the right-lifting convention
\[
    [\Phi,\Psi](b,a)=\int_c[\Phi(a,c),\Psi(b,c)],
\]
the first variable of \(\ll_X(a,b)\) comes from the source variable of
\((\eta_X)_\ast\), while the second variable comes from the source variable of
\(\alpha_X^\ast\). Thus
\[
    \ll_X(a,b)
    =
    [\alpha_X^\ast,(\eta_X)_\ast](a,b)
    =
    \int_I
    [
        \alpha_X^\ast(b,I),
        (\eta_X)_\ast(a,I)
    ].
\]

The universal right-lifting triangle is
\[
\begin{tikzcd}[column sep=huge,row sep=large]
X
    \ar[r,dashed,"\ll_X"]
    \ar[rr,bend right=18,dashed,"{(\eta_X)_\ast}"']
&
X
    \ar[r,dashed,"\alpha_X^\ast"]
&
\operatorname{Idl}_{\mathbb V}(X)
\ar[phantom,from=1-2,to=1-3,"\Downarrow" description]
\end{tikzcd}
\]
and its universal property says that every compatible triangle factors through
\(\ll_X\):
\[
    \operatorname{Dist}_{\mathbb V}(X,X)(\Theta,\ll_X)
    \cong
    \operatorname{Dist}_{\mathbb V}
    \bigl(X,\operatorname{Idl}_{\mathbb V}(X)\bigr)
    \bigl(\alpha_X^\ast\circ\Theta,(\eta_X)_\ast\bigr).
\]
Equivalently, a comparison
\[
    \alpha_X^\ast\circ\Theta\to(\eta_X)_\ast
\]
is the same as a comparison
\[
    \Theta\to\ll_X.
\]

\begin{proposition}[Formula for the fuzzy way-below distributor]
\label{prop:fuzzy-way-below-formula}
For every fuzzy domain \(X\),
\[
    \ll_X(a,b)
    \cong
    \int_{I\in \operatorname{Idl}_{\mathbb V}(X)}
    \left[
        X(b,\alpha_X I),
        \operatorname{Idl}_{\mathbb V}(X)(\eta_X a,I)
    \right].
\]
By enriched Yoneda, this is equivalently
\[
    \ll_X(a,b)
    \cong
    \int_{I\in \operatorname{Idl}_{\mathbb V}(X)}
    \left[
        X(b,\alpha_X I),
        I(a)
    \right].
\]
\end{proposition}

\begin{proof}
This is the right-lifting formula applied to
\[
    \Phi=\alpha_X^\ast,
    \qquad
    \Psi=(\eta_X)_\ast.
\]
Thus
\[
    \ll_X(a,b)
    =
    [\alpha_X^\ast,(\eta_X)_\ast](a,b)
\]
is
\[
    \int_{I\in \operatorname{Idl}_{\mathbb V}(X)}
    [
        \alpha_X^\ast(b,I),
        (\eta_X)_\ast(a,I)
    ].
\]
Using the definitions of \(\alpha_X^\ast\) and \((\eta_X)_\ast\), this becomes
\[
    \int_{I\in \operatorname{Idl}_{\mathbb V}(X)}
    \left[
        X(b,\alpha_X I),
        \operatorname{Idl}_{\mathbb V}(X)(\eta_X a,I)
    \right].
\]
Finally,
\[
    \operatorname{Idl}_{\mathbb V}(X)(\eta_X a,I)
    \cong
    I(a)
\]
by enriched Yoneda.
\end{proof}

\begin{proposition}[Universal property of way-below]
\label{prop:way-below-universal-property}
For every fuzzy distributor
\[
    \Theta:X\nrightarrow X,
\]
there is a natural isomorphism
\[
    \operatorname{Dist}_{\mathbb V}(X,X)(\Theta,\ll_X)
    \cong
    \operatorname{Dist}_{\mathbb V}
    \bigl(X,\operatorname{Idl}_{\mathbb V}(X)\bigr)
    \bigl(\alpha_X^\ast\circ\Theta,(\eta_X)_\ast\bigr).
\]
In the locally posetal reflection this becomes
\[
    \Theta\leq\ll_X
    \quad\Longleftrightarrow\quad
    \alpha_X^\ast\circ\Theta\leq(\eta_X)_\ast.
\]
\end{proposition}

\begin{proof}
This is the universal property of the right lifting
\[
    [\alpha_X^\ast,(\eta_X)_\ast].
\]
\end{proof}

The universal property can be visualized as the adjunction
\[
\begin{tikzcd}[column sep=huge,row sep=large]
\operatorname{Dist}_{\mathbb V}(X,X)
    \ar[r,bend left=35,"{\alpha_X^\ast\circ -}"]
&
\operatorname{Dist}_{\mathbb V}
    \bigl(X,\operatorname{Idl}_{\mathbb V}(X)\bigr)
    \ar[l,bend left=35,"{[\alpha_X^\ast,-]}"]
\end{tikzcd}
\]
evaluated at \((\eta_X)_\ast\).

\begin{proposition}[Way-below implies below]
\label{prop:fuzzy-way-below-implies-below}
For every fuzzy domain \(X\), there is a canonical comparison
\[
    \ll_X(a,b)\to X(a,b).
\]
In the posetal reflection this says
\[
    \ll_X(a,b)\leq X(a,b).
\]
\end{proposition}

\begin{proof}
By Proposition~\ref{prop:fuzzy-way-below-formula}, the end defining
\(\ll_X(a,b)\) maps to its component at the principal ideal \(\eta_X b\):
\[
    \ll_X(a,b)
    \to
    \left[
        X(b,\alpha_X\eta_X b),
        \operatorname{Idl}_{\mathbb V}(X)(\eta_X a,\eta_X b)
    \right].
\]
The strict algebra unit law gives
\[
    \alpha_X\eta_X b=b.
\]
The Yoneda embedding is fully faithful, so
\[
    \operatorname{Idl}_{\mathbb V}(X)(\eta_X a,\eta_X b)
    \cong
    X(a,b).
\]
Thus we have a map
\[
    \ll_X(a,b)
    \to
    [X(b,b),X(a,b)].
\]
The unit map
\[
    I\to X(b,b)
\]
induces, by contravariance of the first variable of the internal hom, a morphism
\[
    [X(b,b),X(a,b)]
    \to
    [I,X(a,b)]
    \cong
    X(a,b).
\]
The composite gives
\[
    \ll_X(a,b)\to X(a,b).
\]
\end{proof}

The proof follows the diagram
\[
\begin{tikzcd}[column sep=huge,row sep=large]
{\ll_X(a,b)}
    \ar[r]
    \ar[rr,bend right=18]
&
{\left[
    X(b,\alpha_X\eta_X b),
    \operatorname{Idl}_{\mathbb V}(X)(\eta_X a,\eta_X b)
\right]}
    \ar[r,"{\alpha_X\eta_X b=b}"]
&
{\left[
    X(b,b),
    X(a,b)
\right]}
    \ar[r]
&
{X(a,b)}.
\end{tikzcd}
\]

\subsection{The posetal reflection}
\label{subsec:fuzzy-way-below-posetal}

For a fuzzy domain over \(\Lambda\),
\[
    \ll_X(a,b)
    =
    \bigmeet_{I\in\operatorname{Idl}_\Lambda(X)}
    \left(
        X(b,\alpha_X I)\multimap I(a)
    \right).
\]
The distributor law for
\[
    \ll_X:X\nrightarrow X
\]
gives fuzzy monotonicity:
\[
    X(a',a)\tensor \ll_X(a,b)\tensor X(b,b')
    \leq
    \ll_X(a',b').
\]
The comparison of Proposition~\ref{prop:fuzzy-way-below-implies-below} becomes
\[
    \ll_X(a,b)\leq X(a,b).
\]
Thus the way-below degree is always bounded above by the ordinary fuzzy comparison
degree.

In the Lawvere case this inequality is in the truth-value order. Since the
truth-value order is the reverse of the usual numerical order, it reads
numerically as
\[
    \ll_X(a,b)\geq d_X(a,b).
\]

\subsection{Examples}
\label{subsec:fuzzy-way-below-examples}

\begin{example}[The crisp case]
\label{ex:fuzzy-way-below-crisp}
For \(\Lambda=2\), the formula becomes
\[
    a\ll_X b
    \quad\Longleftrightarrow\quad
    \forall I\in\operatorname{Idl}(X),\;
    b\leq \bigvee I\Rightarrow a\in I.
\]
This is the ordinary way-below relation on a dcpo.
\end{example}

\begin{example}[G\"odel way-below degrees]
\label{ex:fuzzy-way-below-godel}
For G\"odel truth values,
\[
    \ll_X(a,b)
    =
    \bigmeet_{I}
    \left(
        X(b,\alpha_X I)\multimap_{\mathrm G} I(a)
    \right).
\]
This is the degree to which every ideal whose join lies above \(b\) contains
\(a\).
\end{example}

\begin{example}[{\L}ukasiewicz way-below degrees]
\label{ex:fuzzy-way-below-lukasiewicz}
For {\L}ukasiewicz truth values,
\[
    \ll_X(a,b)
    =
    \bigmeet_{I}
    \min(1,1-X(b,\alpha_X I)+I(a)).
\]
\end{example}

\begin{example}[Product way-below degrees]
\label{ex:fuzzy-way-below-product}
For product truth values,
\[
    \ll_X(a,b)
    =
    \bigmeet_I
    \left(
        X(b,\alpha_X I)\multimap_\Pi I(a)
    \right).
\]
\end{example}

\begin{example}[Lawvere way-below costs]
\label{ex:fuzzy-way-below-lawvere}
For the Lawvere quantale, the meet in the truth-value order is a supremum in the
usual numerical order. Hence
\[
    \ll_X(a,b)
    =
    \sup_{I\in\operatorname{Idl}_\Lambda(X)}
    \bigl(I(a)\ominus X(b,\alpha_X I)\bigr),
\]
where
\[
    v\ominus u=\max(0,v-u)
\]
with the extended-value conventions. Thus \(\ll_X(a,b)\) measures the
worst-case residual cost needed to turn evidence that \(b\) lies below an ideal
colimit into evidence that \(a\) belongs to that ideal.
\end{example}

\section{Continuity as representability of way-below}
\label{sec:continuity-representability-way-below}

We now show that the previous chapter's way-below formula is precisely the
representable form of the right-lifted distributor.

\subsection{The enriched construction}
\label{subsec:continuity-representability-enriched}

\begin{definition}[Continuous fuzzy domain, recalled]
\label{def:continuous-fuzzy-domain-recalled-equipment}
A fuzzy domain \(X\) is \textbf{continuous} if its ideal-algebra map
\[
    \alpha_X:\operatorname{Idl}_{\mathbb V}(X)\to X
\]
has an enriched left adjoint
\[
    \beta_X:X\to\operatorname{Idl}_{\mathbb V}(X),
    \qquad
    \beta_X\dashv\alpha_X.
\]
\end{definition}

The continuity data form the adjoint diagram
\[
\begin{tikzcd}[column sep=huge,row sep=large]
X
    \ar[r,bend left=35,"\beta_X"]
&
\operatorname{Idl}_{\mathbb V}(X)
    \ar[l,bend left=35,"\alpha_X"]
\end{tikzcd}
\]
and hence give natural isomorphisms
\[
    X(b,\alpha_X I)
    \cong
    \operatorname{Idl}_{\mathbb V}(X)(\beta_X b,I).
\]

\begin{theorem}[Continuity represents the way-below distributor]
\label{thm:continuity-represents-way-below}
Let \(X\) be a continuous fuzzy domain. Then
\[
    \ll_X(a,b)
    \cong
    \operatorname{Idl}_{\mathbb V}(X)(\eta_X a,\beta_X b).
\]
Consequently,
\[
    \ll_X(a,b)\cong \beta_X(b)(a).
\]
Equivalently, as distributors,
\[
    \ll_X
    \cong
    \beta_X^\ast\circ(\eta_X)_\ast.
\]
\end{theorem}

\begin{proof}
By Proposition~\ref{prop:fuzzy-way-below-formula},
\[
    \ll_X(a,b)
    \cong
    \int_I
    \left[
        X(b,\alpha_X I),
        \operatorname{Idl}_{\mathbb V}(X)(\eta_X a,I)
    \right].
\]
Since
\[
    \beta_X\dashv\alpha_X,
\]
we have
\[
    X(b,\alpha_X I)
    \cong
    \operatorname{Idl}_{\mathbb V}(X)(\beta_X b,I).
\]
Therefore
\[
    \ll_X(a,b)
    \cong
    \int_I
    \left[
        \operatorname{Idl}_{\mathbb V}(X)(\beta_X b,I),
        \operatorname{Idl}_{\mathbb V}(X)(\eta_X a,I)
    \right].
\]
By enriched Yoneda in \(\operatorname{Idl}_{\mathbb V}(X)\),
\[
    \int_I
    \left[
        \operatorname{Idl}_{\mathbb V}(X)(\beta_X b,I),
        \operatorname{Idl}_{\mathbb V}(X)(\eta_X a,I)
    \right]
    \cong
    \operatorname{Idl}_{\mathbb V}(X)(\eta_X a,\beta_X b).
\]
Applying Yoneda again to the ideal
\[
    \beta_X b:X^{\op}\to\underline{\mathbb V}
\]
gives
\[
    \operatorname{Idl}_{\mathbb V}(X)(\eta_X a,\beta_X b)
    \cong
    \beta_X(b)(a).
\]

For the distributor-level statement, compute
\[
\begin{aligned}
(\beta_X^\ast\circ(\eta_X)_\ast)(a,b)
&=
\int^{I\in\operatorname{Idl}_{\mathbb V}(X)}
\operatorname{Idl}_{\mathbb V}(X)(\eta_Xa,I)
\tensor
\operatorname{Idl}_{\mathbb V}(X)(I,\beta_Xb)
\\
&\cong
\operatorname{Idl}_{\mathbb V}(X)(\eta_Xa,\beta_Xb)
\\
&\cong
\ll_X(a,b),
\end{aligned}
\]
by enriched co-Yoneda and the calculation above.
\end{proof}

Thus the approximant map
\[
    \beta_X:X\to\operatorname{Idl}_{\mathbb V}(X)
\]
represents the right-lifted way-below distributor.

\subsection{The posetal reflection}
\label{subsec:continuity-representability-posetal}

For a continuous fuzzy domain over \(\Lambda\),
\[
    \ll_X(a,b)
    =
    \operatorname{Idl}_\Lambda(X)(\eta_X a,\beta_X b)
    =
    \beta_X(b)(a).
\]
Thus the way-below degree
\[
    \operatorname{wb}_X(a,b)=\beta_X(b)(a)
\]
from the previous chapter is exactly the value of the right-lifted distributor
constructed in this chapter.

If \(X\) is separated, then the reconstruction theorem from the previous chapter
gives
\[
    \alpha_X\beta_X=1_X.
\]
For completeness, we recall the argument. The unit of
\[
    \beta_X\dashv\alpha_X
\]
gives
\[
    1_X\leq \alpha_X\beta_X.
\]
The counit gives
\[
    \beta_X\alpha_X\leq 1_{\operatorname{Idl}_\Lambda(X)}.
\]
Applying this to the principal ideal \(\eta_X b\), and using the algebra law
\[
    \alpha_X\eta_X b=b,
\]
gives
\[
    \beta_X b\leq \eta_X b.
\]
Applying \(\alpha_X\) gives
\[
    \alpha_X\beta_X b
    \leq
    \alpha_X\eta_X b
    =
    b.
\]
Together with
\[
    b\leq \alpha_X\beta_X b,
\]
this gives equality in the induced crisp preorder. Since \(X\) is separated,
\[
    \alpha_X\beta_X b=b.
\]

The reconstruction is captured by
\[
\begin{tikzcd}[column sep=huge,row sep=large]
X
    \ar[r,"\beta_X"]
    \ar[dr,"1_X"']
&
\operatorname{Idl}_\Lambda(X)
    \ar[d,"\alpha_X"]
\\
&
X
\end{tikzcd}
\]
in the separated posetal reflection.

Hence every object \(b\) of a separated continuous fuzzy domain is recovered as
the fuzzy ideal colimit of its approximant ideal:
\[
    b=\alpha_X(\beta_X b).
\]

\subsection{Examples}
\label{subsec:continuity-representability-examples}

\begin{example}[The crisp case]
\label{ex:continuity-representability-crisp}
For \(\Lambda=2\), continuity means that
\[
    \alpha_X:\operatorname{Idl}(X)\to X
\]
has a left adjoint
\[
    \beta_X:X\to\operatorname{Idl}(X).
\]
Then
\[
    \beta_X(b)=\{a\in X\mid a\ll b\},
\]
and
\[
    a\ll b
    \quad\Longleftrightarrow\quad
    a\in \beta_X(b).
\]
The reconstruction equation
\[
    \alpha_X\beta_X(b)=b
\]
says that every element is the directed supremum of its way-below approximants.
\end{example}

\begin{example}[G\"odel continuous fuzzy domains]
\label{ex:continuity-representability-godel}
For G\"odel truth values,
\[
    \operatorname{wb}_X(a,b)
    =
    \beta_X(b)(a)
\]
is the G\"odel-valued degree to which \(a\) approximates \(b\).
\end{example}

\begin{example}[{\L}ukasiewicz continuous fuzzy domains]
\label{ex:continuity-representability-lukasiewicz}
For {\L}ukasiewicz truth values,
\[
    \operatorname{wb}_X(a,b)
    =
    \beta_X(b)(a)
\]
is an MV-valued approximation degree. The right-lifting formula computes it
using the {\L}ukasiewicz residual.
\end{example}

\begin{example}[Product continuous fuzzy domains]
\label{ex:continuity-representability-product}
For product truth values,
\[
    \operatorname{wb}_X(a,b)
    =
    \beta_X(b)(a),
\]
and the comparison between ideal-colimit evidence and ideal-membership evidence
is governed by the product residual.
\end{example}

\begin{example}[Lawvere continuous fuzzy domains]
\label{ex:continuity-representability-lawvere}
For the Lawvere quantale,
\[
    \beta_X(b)(a)
\]
is the cost with which \(a\) approximates \(b\). Continuity says that the
right-lifted approximation cost is represented by the finite-flat weight
\[
    \beta_X(b).
\]
\end{example}

\section{Scott opens through way-below approximants}
\label{sec:scott-opens-through-way-below}

The previous chapter defined fuzzy Scott opens as upper fuzzy predicates
preserving fuzzy ideal joins. The way-below distributor reconstructs Scott opens
from their values on approximants.

\subsection{The enriched construction}
\label{subsec:scott-opens-way-below-enriched}

Let \(X\) be a weak fuzzy domain with structure map
\[
    \alpha_X:\operatorname{Idl}_{\mathbb V}(X)\to X.
\]
A fuzzy Scott open is an upper fuzzy predicate
\[
    U:X\to\underline{\mathbb V}
\]
satisfying
\[
    U(\alpha_X I)
    \cong
    \int^{x\in X}
    I(x)\tensor U(x)
\]
for every fuzzy ideal \(I\).

The right-hand side is the enriched pairing of the lower weight
\[
    I:X^{\op}\to\underline{\mathbb V}
\]
with the upper predicate
\[
    U:X\to\underline{\mathbb V}.
\]
Thus Scott openness says that \(U\) sends the chosen ideal colimit of \(I\) to
the corresponding weighted colimit of its values. The pairing is represented by
\[
\begin{tikzcd}[column sep=huge,row sep=large]
X^{\op}
    \ar[r,"I"]
&
\underline{\mathbb V}
&
X
    \ar[l,"U"']
\\
&
\int^{x\in X} I(x)\tensor U(x)
    \ar[u,dashed]
\end{tikzcd}
\]
and the Scott condition by
\[
\begin{tikzcd}[column sep=huge,row sep=large]
\operatorname{Idl}_{\mathbb V}(X)
    \ar[r,"\alpha_X"]
    \ar[dr,"{I\mapsto\int^{x}I(x)\tensor U(x)}"']
&
X
    \ar[d,"U"]
\\
&
\underline{\mathbb V}.
\end{tikzcd}
\]

If an enriched reconstruction equivalence
\[
    \alpha_X\beta_X\cong 1_X
\]
is available, then the Scott-open condition gives
\[
    U(b)
    \cong
    U(\alpha_X(\beta_X b))
    \cong
    \int^{a\in X}
        \beta_X(b)(a)\tensor U(a)
    \cong
    \int^{a\in X}
        \ll_X(a,b)\tensor U(a).
\]
In the posetal separated case, the reconstruction equality was obtained in the
previous section, so the formula becomes an equality of truth values.

The reconstruction path is
\[
\begin{tikzcd}[column sep=huge,row sep=large]
X
    \ar[r,"\beta_X"]
    \ar[rr,bend left=22,"1_X"]
&
\operatorname{Idl}_{\mathbb V}(X)
    \ar[r,"\alpha_X"]
    \ar[d,"{I\mapsto\int^{a}I(a)\tensor U(a)}"']
&
X
    \ar[d,"U"]
\\
&
\underline{\mathbb V}
    \ar[r,equal]
&
\underline{\mathbb V}.
\end{tikzcd}
\]

\subsection{The posetal reflection}
\label{subsec:scott-opens-way-below-posetal}

In the posetal specialization, a fuzzy Scott open is an upper fuzzy predicate
\[
    U:X\to L
\]
satisfying
\[
    U(\alpha_X I)
    =
    \bigjoin_{x\in X}
    I(x)\tensor U(x)
\]
for every fuzzy ideal \(I\).

\begin{proposition}[Scott opens are reconstructed from way-below approximants]
\label{prop:scott-open-reconstruction-way-below}
Let \(X\) be a separated continuous fuzzy domain over \(\Lambda\), and let
\[
    U:X\to L
\]
be a fuzzy Scott open. Then for every \(b\in X\),
\[
    U(b)
    =
    \bigjoin_{a\in X}
    \ll_X(a,b)\tensor U(a).
\]
\end{proposition}

\begin{proof}
Since \(X\) is separated and continuous,
\[
    b=\alpha_X(\beta_X b).
\]
Since \(U\) is Scott open,
\[
    U(b)
    =
    U(\alpha_X(\beta_X b))
    =
    \bigjoin_{a\in X}
    \beta_X(b)(a)\tensor U(a).
\]
By Theorem~\ref{thm:continuity-represents-way-below},
\[
    \beta_X(b)(a)=\ll_X(a,b).
\]
Substituting gives
\[
    U(b)
    =
    \bigjoin_{a\in X}
    \ll_X(a,b)\tensor U(a).
\]
\end{proof}

Equivalently,
\[
    U(b)
    =
    \bigjoin_{a\in X}
    \beta_X(b)(a)\tensor U(a).
\]

\subsection{Examples}
\label{subsec:scott-opens-way-below-examples}

\begin{example}[The crisp case]
\label{ex:scott-opens-way-below-crisp}
For \(\Lambda=2\), if \(X\) is a continuous dcpo and \(U\subseteq X\) is Scott
open, then
\[
    b\in U
    \quad\Longleftrightarrow\quad
    \exists a\in X,\; a\ll b\text{ and }a\in U.
\]
Thus Scott opens are determined by way-below approximants.
\end{example}

\begin{example}[G\"odel Scott opens]
\label{ex:scott-opens-way-below-godel}
For G\"odel truth values, the reconstruction formula is
\[
    U(b)
    =
    \bigjoin_{a\in X}
    \min(\ll_X(a,b),U(a)).
\]
\end{example}

\begin{example}[{\L}ukasiewicz Scott opens]
\label{ex:scott-opens-way-below-lukasiewicz}
For {\L}ukasiewicz truth values,
\[
    U(b)
    =
    \bigjoin_{a\in X}
    \max(0,\ll_X(a,b)+U(a)-1).
\]
\end{example}

\begin{example}[Product Scott opens]
\label{ex:scott-opens-way-below-product}
For product truth values,
\[
    U(b)
    =
    \bigjoin_{a\in X}
    \ll_X(a,b)U(a).
\]
\end{example}

\begin{example}[Lawvere Scott opens]
\label{ex:scott-opens-way-below-lawvere}
For the Lawvere quantale, the Scott reconstruction formula becomes
\[
    U(b)
    =
    \inf_{a\in X}
    \bigl(\ll_X(a,b)+U(a)\bigr),
\]
written in the usual numerical order. Thus the cost of satisfying the Scott
predicate at \(b\) is obtained by optimizing over approximants \(a\) of \(b\).
\end{example}

\section{Summary}
\label{sec:summary-fuzzy-distributors-way-below}

This chapter constructed the fuzzy distributor equipment and used it to explain
the origin of the fuzzy way-below relation.

The construction starts with the data already present in the fuzzy ideal monad:
\[
    \eta_X:X\to\operatorname{Idl}_{\mathbb V}(X),
    \qquad
    \alpha_X:\operatorname{Idl}_{\mathbb V}(X)\to X.
\]
The first is the principal-ideal embedding, and the second is the algebra map of a
fuzzy domain. Their strict algebra interaction is
\[
\begin{tikzcd}[column sep=huge,row sep=large]
X
    \ar[r,"\eta_X"]
    \ar[dr,"1_X"']
&
\operatorname{Idl}_{\mathbb V}(X)
    \ar[d,"\alpha_X"]
\\
&
X.
\end{tikzcd}
\]

Fuzzy distributors
\[
    \Phi:X\nrightarrow Y
\]
are enriched profunctors
\[
    X^{\op}\tensor Y\to\underline{\mathbb V}.
\]
They compose by coends, admit companions and conjoints for fuzzy functors, and
have right liftings
\[
    [\Phi,\Psi](b,a)
    =
    \int_c
    [\Phi(a,c),\Psi(b,c)].
\]
The distributor equipment has the schematic form
\[
\begin{tikzcd}[column sep=huge,row sep=large]
X
    \ar[d,"f"']
    \ar[r,dashed,"\Phi"]
&
Y
    \ar[d,"g"]
\\
X'
    \ar[r,dashed,"\Psi"']
&
Y'
\ar[phantom,from=1-2,to=2-1,"\Downarrow" description]
\end{tikzcd}
\]
where vertical arrows are fuzzy functors, horizontal arrows are fuzzy
distributors, and cells are distributor transformations.

For a fuzzy domain \(X\), the conjoint of \(\alpha_X\) and the companion of
\(\eta_X\) are distributors
\[
    \alpha_X^\ast,
    \qquad
    (\eta_X)_\ast
    :
    X\nrightarrow \operatorname{Idl}_{\mathbb V}(X).
\]
Their right lifting defines the fuzzy way-below distributor:
\[
    \ll_X
    =
    [\alpha_X^\ast,(\eta_X)_\ast]
    :
    X\nrightarrow X.
\]
The defining triangle is
\[
\begin{tikzcd}[column sep=huge,row sep=large]
X
    \ar[r,dashed,"\ll_X"]
    \ar[rr,bend right=18,dashed,"{(\eta_X)_\ast}"']
&
X
    \ar[r,dashed,"\alpha_X^\ast"]
&
\operatorname{Idl}_{\mathbb V}(X).
\end{tikzcd}
\]
Its objectwise formula is
\[
    \ll_X(a,b)
    \cong
    \int_{I\in \operatorname{Idl}_{\mathbb V}(X)}
    \left[
        X(b,\alpha_X I),
        I(a)
    \right].
\]
In the posetal reflection this becomes
\[
    \ll_X(a,b)
    =
    \bigmeet_{I\in\operatorname{Idl}_\Lambda(X)}
    \left(
        X(b,\alpha_X I)\multimap I(a)
    \right).
\]

Finally, if \(X\) is continuous, with
\[
    \beta_X\dashv\alpha_X,
\]
then the right-lifted way-below distributor is represented by \(\beta_X\):
\[
    \ll_X(a,b)
    \cong
    \operatorname{Idl}_{\mathbb V}(X)(\eta_X a,\beta_X b)
    \cong
    \beta_X(b)(a).
\]
Equivalently, as distributors,
\[
    \ll_X
    \cong
    \beta_X^\ast\circ(\eta_X)_\ast.
\]
Thus the way-below degree introduced in the previous chapter is not extra
structure: it is the value of the right-lifted distributor represented by the
approximant map of a continuous fuzzy domain.